\documentclass[10pt]{article}
\usepackage{amsfonts}
\usepackage{amsmath}
\usepackage{hyperref}
\usepackage{epsfig}
%\usepackage{palatino}
\def\arXiv#1{\href{http://xxx.lanl.gov/abs/#1}{arXiv:#1}}

%%<mathmacros>Some mathematics macros
\def\Dir{\,\,{\raise.15ex\hbox{/}\mkern-12mu D}}
\def\slD{\slash\!\!\!\!D}
\def\slA{\slash\!\!\!\!A}
%\providecommand\slashed[1]{\slash\!\!\!\!{#1}}
\newcommand{\bb}{\begin{equation}}
\newcommand{\bqn}{\begin{eqnarray}}
\newcommand{\eqn}{\end{eqnarray}}
\def\be{\begin{equation}}
\def\ee{\end{equation}}
\def\beq{\begin{eqnarray}}
\def\eeq{\end{eqnarray}}
\newcommand{\pp}{\partial}
\def\g{\gamma}
\def\G{\Gamma}
\def\d{\delta}
\def\D{\Delta}
\def\a{\alpha}
\edef\b{\beta}
\def\l{\lambda}
\def\m{\mu}
\def\n{\nu}
\def\f{\phi}
\def\vf{\varphi}
\providecommand{\tA}{{\tilde A}}
\providecommand{\talpha}{{\tilde\alpha}}
\providecommand{\cplus}{C^{(+)}}
\providecommand{\cminus}{C^{(-)}}
\def\({\left(}
\def\){\right)}
\def\[{\left[}
\def\]{\right]}
%%</mathmacros>

%\usepackage[OT2,OT1]{fontenc}
%\newcommand\cyr{%
%  \renewcommand\rmdefault{wncyr}%
%  \renewcommand\sfdefault{wncyss}%
%  \renewcommand\encodingdefault{OT2}%
%  \normalfont
%  \selectfont}
%\DeclareTextFontCommand{\textcyr}{\cyr}

%\def\refname{{\cyr Literatura}}
%%<somecyrillicmacros>These are some cyrillic macros by me - E.V.
%\def\y{{\cyr\hard}}
%\def\sht{{\cyr shch}}
%\def\Sht{{\cyr SHCH}}
%\def\ik{{\cyr \u i}}
%\def\Ik{{\cyr \u I}}
%\providecommand\hcyr{{\cyr h}}
%\providecommand\hemcyr{{\cyr hem}}
%\providecommand\emcyr{{\cyr em}}
%\providecommand\ya{{\cyr ya}}
%\providecommand\Scyr{{\cyr S}}
%%</somecyrillicmacros>

%\voffset=-2cm
\voffset=-1cm
\hoffset=-2cm
\textheight=21cm
\textwidth=15cm

\begin{document}

\vspace*{1cm}

\begin{center}
{\bf{The UPLNC Compiler: Design and Implementation\\
%\vspace*{.35cm}
%and its implementation for compressing virtually uncompressible data 
}}

\vspace*{1cm}

Evgueniy Vitchev\footnote{e-mail: vitchev@physics.rutgers.edu}
\ \\

\ \\
Department of Physics and Astronomy, Rutgers University \\
Piscataway, NJ 08855-0849, USA
\end{center}


\vspace*{.8cm}

\begin{abstract}
The implementation of the compiler of the UPLNC language is presented
with a full source code listing.
\end{abstract}

\section{Introduction}
\subsection{Why develop another language?}
There were several motivations to develop the UPLNC language. One of
them was to make a programming language which suits better one's
personal tastes, and allows experimenting with ideas. Also, developing
of the compiler was an amusing process 
for the author, who hopes that this work will be instructive to the
reader. From a certain point, one could view such a work as a work of
art, like a poet would write a beautiful poem.
Another reason behind this language would be the beginning of a much
larger project- building an alternative do-it-yourself (DIY) platform,
including programming language, developer tools, operating system
kernel written in that language, full set of system and application
utilities.
If some day such a project becomes implemented, it is the author's hope
that it will remain free, as free speech, without being threatened by
claims of exclusive rights to use, modify and distribute. That's why the
source code in the present work is released under the GNU General Public
License. 

This language is already quite developed- the compiler does support
unlimited level of indirection, multiple dimensional arrays, structures,
methods of structures, and most of the control statements inherent of
C. On the other hand, there are some more or less essential things which
would be desirable to be implemented, e.g. floating point support,
initializers and optimization phase before the code generation. 

An open question is whether this language will evolve as a
general-purpose one like C and C++, or will have some more sophisticated
features as combinatoric control statements(CCS). An lyrical diversion
which gives 
some notion of CCS is a calculation of the form
\be
R=\sum_{1\le i_1\le i_2\le\ldots\le i_k\le N}f(i_1,i_2,\ldots,i_k).
\ee
When $k$ (the number of indices) is not known, the simplest way is to
write some recursive or iterative algorithm which cycles the $k$ indices
in the appropriate way. This can be done also with a standard library
function/template. A non-trivial way to deal with this situation is
post-compile code generation: the program can generate machine code
corresponding to something like
\begin{verbatim}
for(i[k-1]=1;i[k-1]<=N;++i[k-1])
 for(i[k-2]=1;i[k-2]<=i[k-1];++i[k-2])
  for(i[k-3]=1;i[k-3]<=i[k-2];++i[k-3])
   ...
     R+=f(i);
\end{verbatim}
and run that machine code, or it can output the source code to a file,
invoke an external compiler which would produce a shared object, which
then would be loaded and run.
\subsection{Development}
In part, this work was inspired by  Small C compiler \cite{SmCfiles} and
its variant RatC \cite{ABookOnC}. When the author encountered RatC
around one decade ago, he made a rough implementation(or translation) of it
in Turbo Pascal. That project was ultimately abandoned, mainly because
it was a compiler written in a quite high-level language, translating(by
design) a rather primitive language, so a bootstrap procedure was not
possible. Its source code is probably lost. 

Much later, when I saw several other implementations of Small C,
decided that it would be better, instead of making subsets and
dialects, which will probably never be as good as the complete language, to
make one's own language, referred to as UPLNC, which will allow more
choices of its design. Still, I decided to use Small C as the first language to
implement the UPLNC compiler-- to that end I ported a version of Small C
to \texttt{i386-linux} \cite{SmCport}. Those preliminary versions of the
compiler 
were pretty much constrained by the paradigms of the Small C compiler--
no use of structures at all, so that tables were implemented by means of
several arrays, corresponding to the different fields of the missing
structures. Even in the present version one can find similarities of the
syntax parser of UPLNC to the one of Small C, as well as artefacts of
code borrowed from the Small C compiler. Nevertheless, the author
wanted to escape the constraints of Small C, that's why, when the
compiler was already able to produce correct code, the author wrote a
utility for translating Small C to UPLNC (one can see its traces by some
truncated identifiers like \texttt{expressi}). The result was a compiler for
UPLNC, written in the same language, so it could compile itself.

A next step was to replace the implementation of the tables as arrays of
scalar types with arrays or lists of structures. The processing of
expressions was completely rewritten. The old Small C-style approach of
code generation being done while parsing, was replaced by a three-stage
process: parsing and building a tree, generating from that three an
intermediate representation of the program code, producing of assembly
code. It is possible in the future to include an optimization stage
after the second 
stage of this process.

Later, methods of structures were introduced in the language, much
similar to the methods in C++: the method has a default parameter
{\textsf{this}} which points to the instance of the structure, and within
the method the members of the structure are visible.

It is possible to declare variables in different styles:
\begin{verbatim}
var a,b,c:[3]*char;
var [3]*char:d,e,f;
var extern a1,a2:int;
var extern int:a3;
var int extern:a4;
var a5:extern int;
var a6:int extern;
var extern a7:extern int extern;
var extern **int extern: extern a8;
\end{verbatim}

Currently, the only supported target of the compiler is Linux on
\textsf{i386}-compatible processors. The code generated is
link-compatible with C, 
that's why it is possible to link programs with the GNU C library
\textsf{glibc}, which is used by the compiler itself. It is plausible
that with small modifications the compiler could support other operating
systems running on \textsf{i386}-compatible microprocessors. In order to
support other architectures, it is possible to rewrite the code
generating part of the compiler. This was the main reason behind
separating the code generation in a different stage and even in the
different module \texttt{codegen.e}.

\begin{figure}[h]
\begin{center}
%\includegraphics[width=10cm]{cylfig1.eps}
%\epsfxsize=0.3\textwidth
%\epsfbox{tlgraph}
%\\
%\epsfxsize=0.42\textwidth
%\epsfbox{tlgraph2}
%\epsfxsize=0.7\textwidth
%\epsfbox{tlgraph3}
\epsfxsize=0.99\textwidth
\epsfbox{tlgraph4}
\end{center}
\caption{Scheme of \texttt{grph.e}}
\label{lfig1}
\end{figure}

A feature introduced in the compiler was the possibility of redefining
certain syntax elements. An illustrative, albeit not quite readable
example is this program for calculating of primes:
\begin{verbatim}
%%takeoff lcomp "<font>" ;
%%takeoff rcomp "</font>" ;
%%takeoff tfunc "<p>" ;
%%takeoff larg "<arg>" ;
%%takeoff rarg "$</arg>" ;

#define MAXNUM 1000
#define TABSZ  1001

var tab:[TABSZ]int;/*factorizes?*/

<p> doprimes<arg>$</arg>
<font>
  var int:i,k;
  printf<arg>"Calculating the primes>=3...\n"$</arg>;
  for<arg>i=2;i<TABSZ;i++$</arg>
    tab[i]=0;
  for<arg>i=3;i<TABSZ;i=i+2$</arg>
    <font>
      if<arg>tab[i]$</arg>
        continue;
      printf<arg>"%d\n",i$</arg>;
      for<arg>k=i;k<TABSZ;k=k+i$</arg>
        <font>
          tab[k]=1;
        </font>
    </font>
</font>
<p> main<arg>$</arg>
<font>
  doprimes<arg>$</arg>;
</font>
\end{verbatim}
\begin{figure}[h]
\begin{center}
%\includegraphics[width=10cm]{cylfig1.eps}
%\epsfxsize=0.3\textwidth
%\epsfbox{tlgraph}
%\\
%\epsfxsize=0.42\textwidth
%\epsfbox{tlgraph2}
%\epsfxsize=0.7\textwidth
%\epsfbox{tlgraph3}
\epsfxsize=0.99\textwidth
\epsfbox{tlgraph5}
\end{center}
\caption{Scheme of \texttt{langc.e}}
\label{lfig2}
\end{figure}

Another feature which was implemented lately was the code graphing,
namely producing a \verb+gnuplot+ \cite{GnuPlot} input file containing a
graphical 
scheme of the referrals between global symbol names in a particular
module. For example, the scheme produced for some of the modules in
the compiler are shown in  Fig.\ref{lfig1} and Fig.\ref{lfig2}.


The compiler driver was written in Perl \cite{PerlOrg}. It runs the
preprocessor \texttt{lpp1} 
and redirects its output as input of the compiler
\texttt{langc}. Linking is done by using the driver of \texttt{gcc}
\cite{GNUCC}. 

\section{Source}
In this section a full source code listing of the compiler is
presented. An archive of previous versions is in \cite{UPLNChome}.
\subsection{langc.e}
\input{langc.e.inc}
\subsection{codegen.e}
\input{codegen.e.inc}
\subsection{grph.e}
\input{grph.e.inc}
\subsection{tlangc.he}
\input{tlangc.he.inc}
\subsection{codegen.he}
\input{codegen.he.inc}
\subsection{autodyn.e}
\input{autodyn.e.inc}
\subsection{lpp1.e}
\input{lpp1.e.inc}
\subsection{langdrv.pl}
\input{langdrv.pl.inc}

\begin{thebibliography}{99}

\bibitem{SmCfiles}
\href{http://www.cpm.z80.de/small_c.html}
{\texttt{http://www.cpm.z80.de/small\_{}c.html}}

\bibitem{RonCain} Ron Cain, ``A Small C Compiler for the 8080'',
  Dr. Dobb's Journal, May 1980.

\bibitem{ABookOnC}
R.~E.~.Berry, B.~A.~Meekings, ``A Book on C''.

\bibitem{SmCport} Small C port to i386-linux by Evgueniy Vitchev, \\
\href{http://www.physics.rutgers.edu/~vitchev/smallc-i386.html}
{\texttt{http://www.physics.rutgers.edu/$\sim$vitchev/smallc-i386.html}}

\bibitem{UPLNChome} Home page of the UPLNC programming language,\\
\href{http://www.physics.rutgers.edu/~vitchev/thelang/thelang.html}
{\texttt{http://www.physics.rutgers.edu/$\sim$vitchev/thelang/thelang.html}}.

\bibitem{GnuPlot}\href{http://www.gnuplot.info/}{http://www.gnuplot.info/}

\bibitem{GNUCC}\href{http://gcc.gnu.org/}{http://gcc.gnu.org/}

\bibitem{PerlOrg}\href{http://www.perl.org/}{http://www.perl.org/}

\end{thebibliography}

\end{document}

